%! AP Statistics — Unit 7, Lesson 7.3 — Lesson Plan
\documentclass[10pt]{article}
\usepackage{apstats-article}
\usepackage{apstats-boxes}

\newcommand{\UnitNumberName}{Unit 7: Inference for Quantitative Data: Means \quad}
\newcommand{\LessonNumberName}{Lesson 7.3: Justifying a Claim About a Population Mean Based on a Confidence Interval}

\begin{document}
\begin{center}
    {\Large\bfseries \CourseName: \SchoolYear \\
    {\small\bfseries \UnitNumberName \LessonNumberName}}
\end{center}

\begin{tcolorbox}[colback=sky, colframe=navy, boxrule=0.9pt, arc=2mm,
    left=2mm, right=2mm, top=1.5mm, bottom=1.5mm]
    \textbf{Primary Objective:} Students will correctly interpret one-sample $t$-intervals,
    identify relationships between sample size, confidence level, and interval width, and
    use a confidence interval to justify a claim about a population mean. (Big Idea: UNC)
\end{tcolorbox}

\renewcommand{\arraystretch}{1.6}
\begin{skillbox}[Priority Ideas \& Skills]{goldbox}
\begin{minipage}[t]{0.48\linewidth}
    \textbf{AP Skill 4 --- Statistical Argumentation}
    \begin{itemize}
        \item Interpret a confidence interval for a population mean, including
              matched-pairs mean difference (Skill 4.B; UNC-4.S).
        \item Justify a claim based on a confidence interval (Skill 4.D; UNC-4.T).
        \item Identify relationships between sample size, CI width, and margin of
              error (Skill 4.A; UNC-4.U).
    \end{itemize}
\end{minipage}
\hfill
\begin{minipage}[t]{0.48\linewidth}
    \textbf{Key Understandings}
    \begin{itemize}
        \item A 95\% CI does \emph{not} mean ``95\% probability that $\mu$ is in
              the interval'' --- it means the method captures $\mu$ in 95\% of
              repeated samples (UNC-4.S.1).
        \item A correct interpretation references the specific sample, context,
              and population (UNC-4.S.3).
        \item Increasing $n$ decreases width; increasing confidence level increases
              width; the two trade off (UNC-4.U).
        \item To justify a claim: if the claim value is outside the CI, there is
              evidence against it; if it is inside, there is insufficient evidence
              against it (UNC-4.T.1).
    \end{itemize}
\end{minipage}
\end{skillbox}

\begin{skillbox}[{Vocabulary, Concepts, \& Theorems}]{greenbox}
\renewcommand{\arraystretch}{1.3}
\begin{tabularx}{\linewidth}{lX}
    \textbf{Interpretation of a CI} &
        ``We are $C$\% confident that the interval $(\ell, u)$ captures the true mean
        [parameter] for [population].'' (UNC-4.S.2) \\
    \textbf{What $C$\% confident means} &
        In repeated random sampling with the same $n$, approximately $C$\% of all
        CIs constructed this way will capture $\mu$ (UNC-4.S.1) \\
    \textbf{Width of the interval} &
        $\text{Width} = 2 \cdot ME = 2t^* \cdot s/\sqrt{n}$; proportional to
        $1/\sqrt{n}$ (UNC-4.U.2) \\
    \textbf{Effect of confidence level} &
        Higher $C$ $\Rightarrow$ larger $t^*$ $\Rightarrow$ wider interval
        (UNC-4.U.3) \\
    \textbf{Effect of sample size} &
        Larger $n$ $\Rightarrow$ smaller $SE$ $\Rightarrow$ narrower interval;
        to halve the width, quadruple $n$ (UNC-4.U.1) \\
    \textbf{Using CI to justify a claim} &
        If the claimed value of $\mu$ lies \emph{outside} the CI, the data provide
        evidence against that claim (UNC-4.T.1) \\
\end{tabularx}
\end{skillbox}

\begin{skillbox}[Activate Prior Knowledge \& Spiral Review]{sky}
\begin{minipage}[t]{0.48\linewidth}
    \textbf{Prior Knowledge Check (3--4 min)}
    \begin{itemize}
        \item Lesson 7.2: the full four-step $t$-interval procedure --- conditions,
              formula $\bar x \pm t^* \cdot s/\sqrt{n}$, and $df = n - 1$.
        \item Unit 6: correct vs.\ common-error interpretations of a CI for
              a proportion --- the same logic applies to means.
        \item Unit 6: using a CI to justify a claim about $p$ (analogous to 7.3).
    \end{itemize}
\end{minipage}
\hfill
\begin{minipage}[t]{0.48\linewidth}
    \textbf{Connections Forward}
    \begin{itemize}
        \item One-sample $t$-test (complementary tool) $\to$ Lessons 7.4--7.5
        \item Two-sample $t$-interval (difference of means) $\to$ Lesson 7.6
        \item Selecting the right procedure (CI vs.\ test) $\to$ Lesson 7.10
        \item Same interpretation structure applies to all CIs throughout Units 7--9
    \end{itemize}
\end{minipage}
\end{skillbox}

\begin{skillbox}[Hook \& Lesson]{sky}
\begin{minipage}[t]{0.48\linewidth}
    \textbf{Hook (5--7 min)}\\
    Display a student-generated interval from Lesson 7.2:
    ``95\% CI for mean sodium intake: (2615, 3065) mg.''\\[4pt]
    Show three student ``interpretations'' on the board (one correct, two wrong).
    Ask pairs: ``Which is correct? What is wrong with the others?''\\[4pt]
    Reveal the correct form and discuss what the phrase ``95\% confident'' actually
    means --- a property of the \emph{method}, not of this one interval.
\end{minipage}
\hfill
\begin{minipage}[t]{0.48\linewidth}
    \textbf{Lesson Flow (15 min)}
    \begin{enumerate}
        \item Debrief the three interpretation errors: (a) ``there is a 95\%
              chance $\mu$ is in the interval'' (wrong --- $\mu$ is fixed); (b)
              ``95\% of sample means fall in this interval'' (wrong --- this is about
              the parameter, not the statistic); (c) correct form with reference to
              the sample, context, and population.
        \item Introduce the width relationships graphically: show how doubling $n$
              reduces ME by factor $1/\sqrt{2}$ and quadrupling $n$ halves ME.
        \item Claim-justification framework: ``The claimed value $\mu_0 = 2300$ mg
              lies \emph{outside} the interval (2615, 3065), so the data provide
              evidence that the true mean sodium intake exceeds 2300 mg.''
        \item Matched-pairs interpretation: same language, but $\mu$ refers to the
              mean of \emph{differences} in the population.
    \end{enumerate}
\end{minipage}
\end{skillbox}

\begin{skillbox}[Explicit Instruction: Interpretation \& Claim Justification]{sky}
\begin{minipage}[t]{0.48\linewidth}
    \textbf{Correct Interpretation Template}\\
    ``We are $C$\% confident that the interval \blank{2cm} to \blank{2cm}
    captures the true mean \underline{[quantity]} for \underline{[population]}.''\\[6pt]
    \textbf{Common Errors to Correct}
    \begin{itemize}
        \item Saying ``probability'' instead of ``confident'' (implies $\mu$ could
              be in different places --- it is fixed).
        \item Omitting the population or sample context.
        \item Claiming the interval contains 95\% of the data values.
    \end{itemize}
\end{minipage}
\hfill
\begin{minipage}[t]{0.48\linewidth}
    \textbf{Claim Justification Framework}
    \begin{enumerate}
        \item State the claimed value $\mu_0$.
        \item Report the CI you computed.
        \item If $\mu_0$ is \emph{outside} the CI: ``Because $\mu_0$ is not captured
              by the interval, the data provide convincing evidence that the true mean
              differs from $\mu_0$.''
        \item If $\mu_0$ is \emph{inside} the CI: ``Because $\mu_0$ is plausible
              based on the interval, we do not have convincing evidence against the
              claim.''
        \item One-sided claims: check whether the entire CI lies above (or below)
              the claimed value.
    \end{enumerate}
\end{minipage}
\end{skillbox}

\begin{skillbox}[Active Monitoring]{sky}
\begin{minipage}[t]{0.48\linewidth}
    \textbf{Circulate and Check}
    \begin{itemize}
        \item Flag probability language in interpretations; insist on ``confident.''
        \item Ensure the population is named (not just ``the sample'') in the
              interpretation.
        \item Check that claim justifications reference the CI explicitly, not just
              a vague ``the evidence shows.''
    \end{itemize}
\end{minipage}
\hfill
\begin{minipage}[t]{0.48\linewidth}
    \textbf{Cold-Call Prompts}
    \begin{itemize}
        \item ``A classmate says `there is a 95\% chance $\mu$ is between 2615 and
              3065.' What is wrong with that statement?''
        \item ``If we increase the confidence level to 99\%, what happens to the
              width of the interval and to the margin of error?''
        \item ``The claimed value is inside the interval. Does that mean the claim
              is proved true? Explain.''
    \end{itemize}
\end{minipage}
\end{skillbox}

\begin{skillbox}[Group Work \& Differentiation]{redbox}
\begin{minipage}[t]{0.48\linewidth}
    \textbf{Partner Activity (10--12 min)}\\
    Students receive three completed intervals (different contexts, confidence
    levels, and claimed values). For each:
    \begin{enumerate}
        \item Write a correct one-sentence interpretation.
        \item Determine whether the data provide convincing evidence against the
              claimed value. Justify using the CI.
        \item Predict how the width changes if $n$ is doubled or $C$ is raised
              from 90\% to 99\%.
    \end{enumerate}
\end{minipage}
\hfill
\begin{minipage}[t]{0.48\linewidth}
    \textbf{Differentiation}
    \begin{itemize}
        \item \textit{Support (Tier R)}: Sentence frame provided; only items 1--2.
        \item \textit{On-grade (Tier A)}: All three items for all three scenarios.
        \item \textit{Extension (Tier E)}: For one scenario, determine the minimum
              $n$ needed to reduce the margin of error to a specified value. Justify
              algebraically.
        \item \textit{ELL}: Bilingual sentence stems for the interpretation template.
    \end{itemize}
\end{minipage}
\end{skillbox}

\begin{skillbox}[Individual Work \& Assessment]{redbox}
\begin{minipage}[t]{0.48\linewidth}
    \textbf{Exit Ticket (5 min)}
    \begin{enumerate}
        \item A 95\% $t$-interval for the mean commute time of city workers is
              $(24.3,\ 31.7)$ minutes. Write a correct interpretation.
        \item The city claims the mean commute is 22 minutes. Does the interval
              provide convincing evidence against this claim? Justify.
        \item If the sample size were quadrupled (same $\bar x$ and $s$), describe
              the effect on the margin of error.
    \end{enumerate}
\end{minipage}
\hfill
\begin{minipage}[t]{0.48\linewidth}
    \textbf{AP-Style Check}\\
    \textit{A 95\% confidence interval for the mean weight (oz) of a product is
    $(15.8,\ 16.2)$. A manager claims the mean weight is 16.5 oz. Which of the
    following is the best conclusion?}
    \begin{enumerate}[label=(\Alph*)]
        \item The mean weight is definitely not 16.5 oz.
        \item There is convincing evidence that the mean weight is not 16.5 oz.
        \item There is not convincing evidence the mean weight differs from 16.5 oz.
        \item 95\% of products weigh between 15.8 and 16.2 oz.
    \end{enumerate}
    Answer: (B) --- 16.5 is outside the interval.
\end{minipage}
\end{skillbox}

\begin{skillbox}[Reinforcement \& Extension]{goldbox}
\begin{minipage}[t]{0.48\linewidth}
    \textbf{Homework / Practice}
    \begin{itemize}
        \item Three problems: each supplies a completed CI and a context. Students
              write a correct interpretation, assess a stated claim, and answer a
              width/ME question (effect of changing $n$ or $C$).
        \item One AP-style free-response question requiring a full four-step interval
              procedure plus interpretation and claim justification.
    \end{itemize}
\end{minipage}
\hfill
\begin{minipage}[t]{0.48\linewidth}
    \textbf{Extension / Preview}
    \begin{itemize}
        \item \textit{Extension}: A 2004 AP FRQ asks students to relate a one-sided
              CI to a significance test. Work through the problem and explain the
              connection.
        \item \textit{Preview}: Lesson 7.4 introduces the one-sample $t$-test ---
              the significance-testing counterpart to the $t$-interval built in
              Lessons 7.2 and 7.3.
    \end{itemize}
\end{minipage}
\end{skillbox}

\end{document}
